\subsection{Physical model and independently checked output}
The study reuses the multilayer flexibility formulation and graph representation
developed in the accompanying project manuscript, while applying the corrected
\texttt{thermal\_transitions\_v1} and \texttt{load\_shifting\_v1} implementations.
The model coordinates thermal generation, nuclear generation, renewable
curtailment, batteries, reservoir and pumped hydro, demand response, line flows
and external trade. Its objective combines operating costs, thermal start-up
costs, flexibility activation, storage cycling, curtailment, unserved energy
and overgeneration penalties. Power-related terms are multiplied by the time
step in hours; start-up costs are charged per event.

Thermal commitment uses exact transition relations, initial conditions and
operating, start-up and shutdown ramp limits. Demand response is lossless load
shifting with a causal energy backlog and terminal repayment; historical prose
describing exponentially disappearing rebound energy does not apply. Storage
tracks efficiencies, retention, initial conditions and terminal bounds. The
network is a capacitated transport model. It does not impose Kirchhoff's laws,
security constraints or individual commitment for every equipment node.
Equipment counts, graph size and the number of free integer decisions are
therefore reported as different quantities.

The shared validator reconstructs physical balances and costs from the complete
dispatch, independently of solver residuals and reported objectives. It checks
finite values, axes, integrality, transitions, ramps, storage, demand-response
repayment, line and trade limits, and nodal balances. The absolute physical and
integrality tolerances are respectively $10^{-3}$ and $10^{-5}$. A feasible
dispatch may include unserved energy because the model permits this penalized
action; the delivered unserved energy is reported separately.

\subsection{Operational success and reference quality}
For scenario $s$, method $m$ and independently allocated budget $b$, let
$F_{smb}$ denote delivery of a physically valid output and normal termination
of the entire process tree within $b$. The measured interval includes scenario
reading, graph construction, online embeddings, candidate generation, repair,
validation, output writing and closure. Static model and library preload is
excluded and reported separately. A forced stop, a late dispatch or an incumbent
owned by a still-running process does not satisfy $F_{smb}$.

A scenario has one common reference cost $C_s^{\rm ref}$, obtained with a
600-second solver limit and 16 CPU threads. We distinguish a certified optimum,
a validated incumbent and a lower bound. A time-limited incumbent is not an
optimum. Training-label computations use their own resource allocation and are
not included in solver timing comparisons. For $|C_s^{\rm ref}|>\epsilon$, define
\begin{equation}
 g_{smb}=\frac{C_{smb}-C_s^{\rm ref}}{|C_s^{\rm ref}|},\qquad
 Y_{smb}(\tau)=F_{smb}\,\ind\{g_{smb}\leq\tau\},\quad \tau=0.05.
 \label{eq:economic}
\end{equation}
Negative gaps are retained: a new method can improve an incumbent reference.
For $|C_s^{\rm ref}|\leq\epsilon=1$ euro, the relative gap is undefined and the
quality condition instead requires $C_{smb}-C_s^{\rm ref}\leq1$ euro. This is an
explicit absolute-tolerance convention, not a relative percentage. Sensitivity
outcomes use $\tau\in\{0.01,0.10\}$, feasibility alone, costs and unserved energy.
If a valid reference is missing, the economic outcome is unknown. An unexecuted
trial is also unknown; neither is encoded as an observed failure.

\subsection{Provenance and statistical units}
The main cohort comprises 600 scenarios, 100 in each of six generator families.
The three persona families and the low, medium and high criticality families
are six distributions, not six ordered severity levels. Four methods are
evaluated at budgets $1,2,5,10,30,60,120$ seconds, giving 16,800 planned primary
trials. Two extra repetitions on 30 preselected cases add 1,680 trials.
Each budget is a new execution; the grid is not a single solver trajectory or
an exact time-to-first-solution observation.

The reserve amendment supersedes identities for affected short budgets while
retaining unchanged 120-second identities, including failures. Analyses reconstruct
the amended schedule rather than enumerate every historical result file.
Source, model, scenario and reference hashes are frozen. Historical files are
read without invoking routines that repair or rewrite checkpoints.
All rows from a scenario, across methods, budgets and repetitions, remain in
one statistical group. A controlled variant shares the group of its base.
Preprocessing is fitted within training groups. The primary coverage analysis
uses the first repetition; repeated trials describe variability separately.
Confidence intervals resample independent scenario groups, not individual rows.
